\clearpage
\phantomsection
\chapter*{摘\quad 要} 
\addcontentsline{toc}{chapter}{摘要}

随着智能制造与工业4.0的深入推进，构建高保真、高实时性的物理实体虚拟副本已成为工业数字孪生技术落地的核心诉求。在这一背景下，虚实映射作为连接物理世界与信息世界的关键使能技术，其底层三维视觉重建底座的质量直接决定了孪生系统的可用性。近年来，以三维高斯溅射（3D Gaussian Splatting, 3DGS）为代表的显式神经渲染技术，凭借其高保真的画面渲染与光栅化效率，为新视角合成带来了新的解决方案。然而，当直接将其应用于包含高反光金属部件、大面积纯色弱纹理外壳以及复杂设备遮挡的数字孪生场景时，传统3DGS暴露出一定的局限性：一方面，纯视觉光度约束在复杂材质上易引发几何拓扑歧义与“浮空伪影”；另一方面，离散高斯基元的无序生长带来了较高的参数冗余，导致显存占用过大，难以满足端侧设备与Web浏览器的大规模部署需求。针对上述挑战，本文以实现面向类工业孪生场景的高效、高保真虚实映射为目标，围绕3DGS的空间几何约束增强与模型轻量化压缩展开了系统性研究。本文的主要研究工作与创新成果如下：

（1）面向类工业场景的多模态自适应几何约束重建算法：针对工业设备表面的各向异性反射与无纹理区域导致的重建病态反问题，本文引入了多模态几何约束。首先，通过空间对齐，将LiDAR获取的物理尺度稠密点云与单目视觉大模型推断的相对深度场进行跨模态配准，为3DGS注入稳健的三维物理初始骨架。其次，设计了基于深度局部高频特征感知的自适应置信度掩膜机制，动态划分平坦信任区与细节豁免区。最后，构建了联合Huber深度惩罚与一阶法向平滑正则化的损失函数，在掩膜的动态调度下，促使高斯基元贴合真实的物理表面。实验表明，该算法有效抑制了金属表面的浮空噪点，改善了纯色机柜的几何撕裂，在三维几何误差（如倒角距离）降低约12\%的同时，较好地保留了机加工纹理的高频边缘。

（2）面向典型孪生场景的3DGS轻量化表达与端侧渲染优化算法：针对3DGS参数体量引发的算力与显存瓶颈，本文从数据重构与渲染管线双管齐下。在几何提纯层面，构建了融合透射率、全局视图可见性以及空间尺度的多维度重要性评估模型，有效剥离了隐藏在不可见区域的冗余“幽灵高斯”。在属性降维层面，引入矢量量化（VQ）技术将存储开销较大的高维球谐函数（SH）系数映射为离散码本，并结合直通估计器（STE）完成量化感知微调（QAT），缓解了硬量化带来的色彩断层。此外，重构了适配Web前端的渲染管线，利用WebGPU计算着色器实现显存内解码，并结合移动端TBDR架构的局部排序优化，打通了轻量化模型在普通终端上的流畅漫游路径。

（3）融合改进3DGS的算法管线构建与系统级综合验证：本文将“多模态几何约束”与“轻量化剪枝量化”在训练管线中进行深度融合，设计了分阶段的渐进式训练策略。研究表明，两者间具有良好的协同效应，稳固的物理几何底座为剪枝操作提供了可靠的结构引导。在开源基准数据集与自主采集的复杂机房孪生数据集上的综合测试表明，本文融合算法在提升图像保真度与几何精度的前提下，有效缓解了系统的计算与存储压力。与传统方法相比，场景模型的物理存储体积减少了约12\%（降至约1010 MB），峰值渲染显存降低至10.2 GB左右，并在常规硬件环境下将渲染帧率提升至42 FPS。该融合算法较好地平衡了渲染质量与系统开销，为数字孪生系统的实际工程部署提供了可行方案。

综上所述，本文的研究成果有效改善了纯视觉3DGS在类工业场景中的材质泛化问题与端侧算力限制，为制造企业构建低算力门槛、高帧率、较好保真度的数字孪生三维视觉底座提供了一套切实可行的轻量化综合解决方案，具有一定的理论意义与工程应用前景。

\vspace{1em}
\noindent\textbf{关键词：} 3D Gaussian Splatting工业场景重建；虚实融合；数字孪生